\chapter{Pulse Oximetry}

\section*{What Is This Test?}
Pulse oximetry is a quick, noninvasive test that estimates peripheral oxygen saturation (SpO$_2$) and displays pulse rate using a light-based sensor, usually on a fingertip. It helps identify low blood oxygen and monitor change over time, but it does not replace clinical assessment or arterial blood-gas testing when precise oxygenation or ventilation data are needed.

\section*{Alternative Names}
Pulse Ox; Oxygen Saturation Test; SpO$_2$ Test; Finger Pulse Oximetry; Oxygen Monitor

\section*{Why It Is Ordered}
Assessment of shortness of breath or suspected hypoxemia; monitoring of lung disease, heart failure, sleep-related breathing problems, or oxygen therapy; and checks in outpatient, emergency, inpatient, procedural, and home-monitoring settings when clinically appropriate.

\section*{How This Test Is Performed}
A clip or adhesive sensor is placed on a finger, toe, or earlobe. Red and infrared light passing through the tissue is analyzed to estimate oxygen saturation and pulse rate. A stable reading requires adequate circulation and minimal movement.

\section*{Preparation, Risks, and Limitations}
Remove dark nail polish when possible and keep the hand warm and still. Poor circulation, motion, skin pigmentation, nail products, bright light, anemia, carbon-monoxide exposure, and device limitations can affect accuracy. A reading should be interpreted with symptoms and the clinical situation; urgent symptoms should not be dismissed because a home device displays a reassuring number.

\section*{References}
\begin{enumerate}
  \item MedlinePlus. Pulse Oximetry.
  \item U.S. Food and Drug Administration. Pulse Oximeter Accuracy and Limitations.
\end{enumerate}

\clearpage
\section*{Understanding Results and Follow-up}
The laboratory report should identify the specimen, method, units, reference interval or decision limit, and whether the result is positive, negative, abnormal, or inconclusive. Reference intervals vary with age, sex, pregnancy, specimen type, assay, and laboratory; a result outside a range is not automatically a diagnosis, and a result within a range does not always exclude disease. Discuss unexpected results with the ordering clinician, who can decide whether repeat or confirmatory testing, treatment, or referral is appropriate. Do not change medicines or delay urgent care based on an isolated result.


\section*{Regulatory Status and Authorization Date}
This chapter describes a test category, not a specific commercial product. FDA, CDSCO, EU IVDR, and MHRA approval or registration dates therefore cannot be assigned without the assay manufacturer, product name, instrument, and intended use. For laboratory-developed tests, an individual agency approval date may not exist. See the Preface for the interpretation of regulatory status.

\clearpage
